% !TeX program = xelatex
%
% This starter is original content written for latextemplates.com; it is not
% the upstream repository's own example deck, which is separately licensed
% CC-BY-NC-SA (non-commercial) and not reused here. Only the theme .sty files
% and the SJTU visual-identity assets (Apache-2.0 / SJTU VI, per LICENSE) are
% vendored from github.com/sjtug/SJTUBeamer.
\documentclass[xcolor=table,aspectratio=169]{ctexbeamer}
\usetheme{sjtubeamer}

\usepackage{booktabs}

\title{学术报告模板}
\subtitle{基于 SJTUBeamer 主题的中英文演示}
\author{王小明}
\institute{上海交通大学 计算机科学与工程系}
\date{\today}

\begin{document}

\begin{frame}
  \titlepage
\end{frame}

\begin{frame}{目录}
  \tableofcontents
\end{frame}

\section{研究背景}

\begin{frame}{为什么要做这项研究}
  \begin{itemize}
    \item 现有方法在大规模场景下计算开销过高
    \item 缺乏统一的评测基准，结果难以横向比较
    \item 本工作提出一套轻量级、可复现的评测框架
  \end{itemize}
\end{frame}

\section{方法}

\begin{frame}{总体流程}
  \begin{columns}[T]
    \begin{column}{0.5\textwidth}
      \textbf{数据准备}
      \begin{itemize}
        \item 采集 12 类真实场景数据
        \item 统一标注格式
        \item 划分训练 / 验证 / 测试集
      \end{itemize}
    \end{column}
    \begin{column}{0.5\textwidth}
      \textbf{Method Overview}
      \begin{itemize}
        \item Feature extraction
        \item Lightweight scoring model
        \item Reproducible evaluation harness
      \end{itemize}
    \end{column}
  \end{columns}
\end{frame}

\begin{frame}{实验结果}
  \begin{table}
    \centering
    \begin{tabular}{lrr}
      \toprule
      方法 (Method) & 准确率 (\%) & 耗时 (ms) \\
      \midrule
      基线 (Baseline) & 71.2 & 340 \\
      本文方法 (Ours) & 89.6 & 118 \\
      \bottomrule
    \end{tabular}
    \caption{在测试集上的性能对比}
  \end{table}
\end{frame}

\section{讨论}

\begin{frame}{关键发现}
  \begin{block}{主要结论}
    在保持较低计算开销的前提下，本文方法将准确率提升了 18.4 个百分点。
  \end{block}
  \begin{itemize}
    \item 在小样本场景下依然保持稳定表现
    \item 对噪声标注具有一定鲁棒性
  \end{itemize}
\end{frame}

\section{总结}

\begin{frame}{总结与展望}
  \begin{itemize}
    \item 提出了一套轻量级评测框架，兼顾准确率与效率
    \item 后续工作将扩展至多模态场景
    \item 代码与数据集将在论文录用后开源
  \end{itemize}
\end{frame}

\begin{frame}{谢谢 / Thank You}
  \centering
  \Large 欢迎提问 \\
  \normalsize wang.xiaoming@sjtu.edu.cn
\end{frame}

\end{document}
